% Vietnamese translation for OI Public Library
% Source-File: IOI中国国家候选队论文/2009/梅诗珂/信息学竞赛中概率问题求解初探.ppt
\documentclass[11pt]{article}
\usepackage{oipl-vn}

\title{Bước vào thế giới xác suất: Bàn sơ về cách giải bài toán xác suất trong thi Tin học\\{\large Ghi chú theo slide}}
\author{Mei Shike\\Trường Trung học số 1 Hợp Phì, An Huy}
\date{2009}

\begin{document}
\maketitle

\origtitle{Solving probability problems in informatics contests}
\sourcepath{IOI-China-Candidate-Team-Papers/2009/Mei-Shike/probability-problems-slides.ppt}

\section{Kiến thức nền}

Slide nhắc lại không gian mẫu, sự kiện, biến ngẫu nhiên, kỳ vọng và xác suất
liên tục. Mục tiêu không phải phát triển lý thuyết xác suất đầy đủ, mà là chỉ
ra cách biến mô hình ngẫu nhiên thành công thức tính được bằng quy hoạch động,
truy hồi hoặc tích phân.

\section{Hai nhóm ví dụ}

\begin{itemize}
  \item Biến ngẫu nhiên rời rạc: \textit{LastMarble} và \textit{Randomness}
  dùng trạng thái hữu hạn, phân phối xác suất và kỳ vọng theo từng bước.
  \item Biến ngẫu nhiên liên tục: \textit{RNG} và \textit{Random Shooting}
  cần mô tả miền hình học hoặc hàm mật độ, rồi đổi bài toán thành tính diện
  tích, thể tích hoặc tích phân.
\end{itemize}

\section{Kinh nghiệm mô hình hóa}

Trước khi viết công thức, cần xác định biến ngẫu nhiên cần hỏi và điều kiện
độc lập giữa các phép thử. Với bài rời rạc, trạng thái nên đủ để mô tả phân
phối tương lai. Với bài liên tục, cần chọn hệ tọa độ và miền tích phân sao cho
ràng buộc trở nên đơn giản.

\section{Ghi chú}

Bản slide là bản dẫn đường cho bài viết đầy đủ. Khi đối chiếu, hãy chú ý cách
mỗi ví dụ chọn đại lượng ngẫu nhiên: chọn đúng đại lượng thường làm công thức
ngắn hơn đáng kể.

\end{document}
